% =====================================================================
%  R. Nielsen, KARAKTER OG VILLIE
%  In: For Ide og Virkelighed. Et Tidsskrift, udgivet af R. Nielsen, Bj. Bjørnson og Rud. Schmidt.  Aargang 1872. Første Bind (Bind 6).
%  Kjøbenhavn: Otto B. Wroblewskys Forlag; G. S. Wibes Bogtrykkeri.  Bind 6, pp. 489-517.  Antiqua.
%
%  TRANSCRIPTION -- GENERATED; DO NOT EDIT BY HAND.  Rasmus Nielsen's eight essays in the volumes 1870-1873 are built by ONE
%  pipeline (ebook-method/for-idee-1870-73/) into a combined working file, which split7.py cuts into seven transcriptions
%  (Aandsvirkelighed and its sequel, Videre om Aandsvirkelighed, are one: Hans, 2026-10-13).
%  Rebuild from the repo root:   sh ebook-method/for-idee-1870-73/rebuild.sh
%
%  STATUS: COMPLETE (2026-10-13): the volume's title page and the whole essay, Bind 6, pp. 489-517, with its notes.
%
%  EDITION.  The only printing.  Text of record: Nasjonalbiblioteket's scan (nb.no) of the volume, ~/bibliotek/Nielsen, Rasmus/1872-1-for-idee-og-virkelighed.pdf (NB item 26d6e6e9175280a596199c9e10bff83c) (SCANS.tsv),
%  page images fetched by nbfetch.py, 300 dpi colour, NO text layer; printed page N = PDF page N + 4.  No e-book, no
%  Wikisource text, no reprint, no earlier transcription found (checked 2026-10-13: the repo, the web, the Internet Archive,
%  da.wikisource; HathiTrust and the Google Books API could not be queried that day; Google Books holds at least the 1873
%  volume, nielsen/darwinismen's source -- a second copy, not used).
%
%  ACCURACY (2026-10-13; one collation for all eight texts, 231 pp., ~58,000 words).  The DRAFT is NB's own OCR (ALTO, ABBYY
%  FineReader 12), laid over the page images (combine.py) and turned into LaTeX by draft.py; it was collated word by word
%  with an independent tesseract reading of the same images (collate/check.py, witness 'ocr'): 1556 disagreements.
%  check.py's own rules settled 78.  Candidate rules were first scored BY SCRIPT on a text already read at the image
%  (calibrate.py: Et Synspunkt for Darwinismen, same volume as Verdensanskuelse), then measured on a BLIND PILOT of 152 crops
%  (reader A, 52/52 planted controls right); 17 rules passed (NB's glyph faults: o for ø, D/O/U for Ø, i for l, li for h,
%  l for I, opening » read «, words run together or broken apart; tesseract's: aa read an/ag/å, added accents; the
%  printer's space before a mark, house rule) and settled 772; the rest were read at the image BLIND (562 crops on 94
%  sheets, four readers, 188/188 planted controls right), ~160 decided in the main session from the readers' lines and ~40
%  at 300-dpi zooms.  STRUCTURE read at the image (reader K): every footnote whole (30; collated where they stand, then moved
%  to their marks, then compared word by word with the reading: notecmp.py, 9 differences settled at the image), the
%  headings, the section numerals, the essays' last lines and signatures, the three Greek words, the volume title pages.
%  LETTERSPACING (the print's emphasis, \emph): candidate lines by letter gaps and box width, calibrated on the words the
%  word readers had marked spaced (spaced.py): 361 lines read blind (61/62 roman controls), 25 doubtful extents read at a
%  zoom.  SWEEPS after OPEN 0: words run together that both OCRs share (lower-case + capital), the unknown-word sweep (678
%  words, ~20 suspects, 14 read at the image: 4 OCR faults, 7 misprints kept), --moved, the dropped-word list, the
%  W-shortens list of the blind decisions.
%  SAMPLED PAGE CHECK (pagecheck.py, seed 202610137): one page of each essay or part, pp. 2289, 2400, 3305, 4308, 5033,
%  5208, 6514, 8011 of the working numbering (Bind x 1000 + page), read whole and blind, 1965 words: 0 word errors,
%  1 punctuation (a compound's hyphen lost at a page turn: „Ny-|Holland“) and 1 letterspacing („to be or not to be“)
%  error, both mended.  95 % upper bound for word errors: 1.5 per 1000 words.
%  What this cannot see: a misprint both OCRs normalise into a real word.
%  The ø rule (Hans 2026-10-08): o-forms that are no word are set with ø; „stode“ (= stood) is a word and stays.
%  PRINTED AS IS: nothing recorded.
%  Not transcribed: running heads, folios, sheet signatures („For Ide og Virkelighed. 1870. I. 20“), library stamps; the
%  large initial of each essay is set as a plain letter.  Quotation marks as printed (» « in the text, sometimes a closing
%  »; „ “ in some notes).  The notes are numbered ¹), ²) per page, as printed.  Greek as printed (textalpha).
%
%  TOOLS (ebook-method/for-idee-1870-73/): nbalto.py, combine.py, alto.py, ocr.py, draft.py, rules.py, calibrate.py,
%  sheets.py + reads.py + READER-BRIEF.md, spaced.py + spaced.tsv, structsheets.py, notecmp.py, unknown.py, zoom.py,
%  pagecheck.py + PAGECHECK-BRIEF.md, patches.py (PATCHES, PRENOTES, NOTES, POST, MISPRINTS), decisions.tsv, split7.py,
%  rebuild.sh.
% =====================================================================
\documentclass[12pt,a4paper,openany]{book}
\usepackage[utf8]{inputenc}
\usepackage[T1]{fontenc}
\usepackage{libertinus}
\usepackage{libertinust1math}
\usepackage{textalpha}
\usepackage[danish]{babel}
\usepackage[protrusion=true,expansion=true]{microtype}
\usepackage{setspace}
\usepackage[margin=1.3in]{geometry}
\usepackage{fancyhdr}
\usepackage{hyperref}
\hypersetup{pdftitle={Karakter og Villie}, pdfauthor={Rasmus Nielsen}, hidelinks}
\pagestyle{fancy}
\fancyhf{}
\fancyhead[LE,RO]{\thepage}
\fancyhead[C]{\textit{R. Nielsen: Karakter og Villie (1872)}}
\setstretch{1.3}
\usepackage{marginnote}
\renewcommand{\thefootnote}{\arabic{footnote})}   % the print numbers each page's notes ¹), ²)
\newcommand{\opage}[1]{\marginnote{\footnotesize\textit{[#1]}}}
% the essay's heading: title (and subtitle), centred, a short rule; the print's large initial is not reproduced
\newcommand{\essayhead}[2]{\begin{center}{\large #1}\ifx&#2&\else\\[0.3em]{\normalsize #2}\fi\\[0.5em]\rule{3em}{0.4pt}\end{center}\medskip\noindent}
\newcommand{\signature}[1]{\par\smallskip\hfill #1\hspace*{2em}\par}

\begin{document}

% --- front matter: the title page of Bind 6 (PDF 1 of 1872-1-for-idee-og-virkelighed.pdf), read at the image (reader K, m1);
% short wavy rules are set as plain rules; the library stamps are not transcribed.
\begin{titlepage}
\begin{center}
\vspace*{3em}
{\Huge For Ide og Virkelighed.}\\[1em]
\rule{8em}{0.4pt}\\[1em]
{\large Et Tidsskrift,}\\[0.8em]
{\small udgivet}\\[0.4em]
{\footnotesize af}\\[0.6em]
{\large R. Nielsen,}\\[0.4em]
{\large Bj. Bjørnson og Rud. Schmidt.}\\[1em]
\rule{8em}{0.4pt}\\[0.6em]
{\large Aargang 1872. Første Bind.}\\[0.6em]
\rule{8em}{0.4pt}\\[4em]
Kjøbenhavn.\\[0.4em]
Otto B. Wroblewskys Forlag.\\[1.5em]
{\small G. S. Wibes Bogtrykkeri.}
\end{center}
\end{titlepage}

% ===== TEXT: pp. 489-517 =====
\essayhead{Karakter og Villie.}{}
% --- p. 489 ---
\setcounter{footnote}{0}%
\opage{489}Zoologisk betragtet, staaer Mennnesket paa det øverste
Trin i Dyrenes Række og indtager sin Hædersplads foran
Aber og Flagermus. Medens ældre Naturkyndige næsten
over Evne anstrengte sig for at gjøre Kløften mellem
Mennesket og de Umælende dyb og bred, har den nyere
Videnskab med Forkjærlighed paaviist og indskærpet det
naturlige Slægtskab mellem Dyr og Mennesker. Vistnok
kan det ikke falde nogen fornuftig Naturforsker ind at
nægte Mennesket fortrinlig Begavelse i Sammenligning
med de øvrige Dyr, men Spørgsmaalet er, hvorledes det
i Grunden staaer til med denne Begavelse, om den tyder
paa noget Oprindeligt eller blot er en Gradforskjel. At
et Menneske har klarere Bevidsthed og mere omfattende
Forstand end en Hund, en Hest, en Ræv, en Elefant, for
ikke at tale om Flagermus, er hævet over al Tvivl; men
sammenligne vi nu en Ræv, der efter Alles Mening har,
hvad man i dyrisk Forstand maa kalde et godt Hoved,
med en Østers, der bogstavelig er hovedløs, saa er her
igjen en betydelig Afstand. Og dog underholder Naturen
Myriader af Individer, som, hvad Intelligensen angaaer,
staae paa lavere Trin end en Østers. Antage vi nu, at
% --- p. 490 ---
\setcounter{footnote}{0}%
\opage{490}alle levende Væsener paa Jordkloden oprindelig ere af
een Slægt, at Organismerne fra Begyndelsen af omtrent
have lignet hinanden, saa at de typiske Forskjelligheder
først efterhaanden i Tidernes Løb ere blevne udviklede;
forestille vi os, at der for Millioner Gange Millioner
Aar siden har været en Periode, da Jordens ypperste
Skabning var en Fisk, en underlig pantseret, tvetydig
Fisk; at Fiskeavlen i Løbet af Millioner Aar tog saa mærkværdige
Opsving, at den tvetydige Fisk, da den omsider
døde, efterlod sig en talrig Familie af Fiskeøgler, Flyveøgler,
Skildpadder, Iguanodonter og Krokodiler, der ligeledes
vare tvetydige; at der saa igjen efter Forløb af
Millioner Gange Billioner Aar indtraadte en Periode, da
de gamle Firbeen avlede nye Firbeen, de Firbenede, vi
kalde Pattedyr; at der blandt disse Firbenede herskede en
saa alsidig Stræben, at nogle, der stræbte opad, bleve til
Fugle, medens andre, der stræbte nedad, bleve til Mulvarpe;
at der til allersidst fremkom et Moderdyr, saa
dobbelt frugtsommeligt, at dets Unger endnu før Fødselen
gik i to Retninger, saa at nogle fik ingen Hænder
og bleve til Marsvin, men andre fik ingen Hale og bleve
til Mennesker: saa have vi i løse Træk et Billede paa
den gradvise Udvikling.

Det er i nærværende Sammenhæng ingenlunde vor
Agt at gaae i Rette med de Naturkyndige om Slægters
og Arters trinvise Dannelse i Jordklodens Udviklingshistorie;
vi holde paa vor gamle Grundsætning, at de
»hedenske, opblæste Videnskaber« (!), uafhængige af »Dogmatiken«,
maae have fuld Ret til at stille sig deres egne
Opgaver og forfølge deres eget Maal efter egen Rettesnor.
Skulde Fysikerne føle Trang til en høiere Rettesnor,
% --- p. 491 ---
\setcounter{footnote}{0}%
\opage{491}skulde Videbegjærligheden friste dem til at gjøre dogmatiske
Indblik i Edens Have, for at nyde Beskuelsen af de
to »absolut uensartede Træer« og see, hvordan Adam tog
sig ud, før Rovdyrene endnu gik paa Rov: da beklager
Filosofien, at den ikke har objektive Vink at tilbyde; de
Videbegjærlige faae i et saa vigtigt Granskningens Anliggende
at søge Raad hos Bisper og Theologer.

Saa glædeligt det end er, at Kjendskab til Naturlærens
vigtigste Sandheder udbredes og optages i den almindelige
Dannelse, saa beklagelig er igjen den Nutidens
Overtro, at Forstaaelsen af det indvortes, sjælelige Liv
skal undergives en naturvidenskabelig Maalestok, at Religion
skal begrundes paa Zoologi, og Samvittigheden udledes
af Fysiken. Det skal atter og atter indskærpes, at
Videnskaben selv er uskyldig i denne Misforstaaelse. Den
sande Videnskab holder sig strengt til Logikens Lov, at
enhver Slutning nøiagtig maa svare til sine Præmisser.
Heraf følger, at Præmisser, der udelukkende tilhøre Naturlæren,
umulig kunne lede til Slutninger, hvis Præmisser
tilhøre Aandslæren. Aandslæren tager sit Udgangspunkt
i Menneskets Indre, dets Grundbevidsthed, dets væsenlige
Selv; men Menneskets indre, væsenlige Selv er lige saa
lidt en Gjenstand for Zoologien som for Geologien. Det
skal indskærpes, at den moderne Aandsforvirring, hvis
Trævlerødder grene sig i en uhyggelig Blanding af overfladiske
Betragtninger og falske Slutninger, hidrører fra
mislig Omgang med objektiv Viden; og saa har den fra
første Færd af sin Grund i det gamle theologiske Bedrag,
som kun skiftede Skin, da man opgav Pavens Autoritet
for at give sig ind under Videnskabens.

Den indvortes Verden har sine Kjendsgjerninger lige
% --- p. 492 ---
\setcounter{footnote}{0}%
\opage{492}saa fuldt som den udvortes. Det er en klar, uomstødelig
Kjendsgjerning, at Mennesket blandt alle levende Skabninger
paa Jorden er den eneste, der har en Dom om
sig selv, den eneste, der kan give sig et Navn og sige ved
sig selv: jeg er et Menneske. I denne Dom ligger, at
Mennesket er et Selv, et Jeg. Med dette Ord: Jeg, er
der sat en uendelig Kløft imellem det klogeste Dyr og
det dummeste Menneske. Vel kan et klogt Dyr under
visse Omstændigheder udvise større Kløgt end et tænksomt
Menneske, og i naturlig Forstand endog siges at
staae langt over en Idiot; men det Blændende i den Art
Sammenligninger er let at hæve. Naar man vil undersøge
Forskjellen mellem Dyr og Menneske, er det dog
virkelig en bagvendt Fremgangsmaade at stille et Dyr
med sunde og et Menneske med syge Organer sammen
eller betragte den Sikkerhed, Dyret har i sit Instinkt,
som ensartet med den Forvisning, Mennesket først maa
tilkæmpe sig ved Eftertanke og mangesidig Overveielse.
Nei, Spørgsmaalet lyder paa noget Andet: hvad er Dyret
og hvad er Mennesket, naar vi see paa de oprindelige
Anlæg? Dyret er en Sjæl uden Aand, Menneskesjælen
er væsenlig Aand; fra dette Synspunkt er Naturgrændsen
tydelig, utvetydig; Dyret er Sandsebevidsthed, Menneskets
Væsen er Selvbevidsthed: dette Grændseskæl staaer fast,
saalænge Dyret er Dyr, og Mennesket Menneske.

Men, siger man og beraaber sig paa Geologien, hele
Dyreriget befinder sig jo i en gradvis Udvikling, en stille,
langsomt opadskridende Bevægelse; om Millioner Gange
Millioner Aar ville Dyrene stifte Stater, Faarene holde
Rigsraad, Hestene studere Geologi, Hønsene emancipere
sig, og Hanen galende sige Jeg til sig selv. Det er en
% --- p. 493 ---
\setcounter{footnote}{0}%
\opage{493}mørk Tale, der ikke tør forstaaes bogstavelig. For det
Første har Geologien ingen Forjættelse om, at de nulevende
Dyreformer ville i en saa fjern Fremtid endnu
være til paa Jorden; den leder os tvertimod til den Formodning,
at de Arter og Slægter, vi kjende, maae, naar
det store Vendepunkt indtræder, alle være uddøde. Det
bliver da ganske vist ikke disse Dyr, der naae op til det
Menneskelige, men andre Dyr, andre Heste, andre Høns,
Heste med Aand, Høns med Selvbevidsthed! For det
Andet er Hypothesen om den gradvise Overgang fra Dyr
til Mennesker, hvad dens Rimelighed angaaer, saa rund som
en Firkant. Thi lad os antage, at der kom en Tid, hvis
det ikke strax kan skee om Millioner Gange Millioner,
saa dog om Billioner Gange Billioner Aar --- Geologien har
jo Tiden for sig --- da Dyrene virkelig fik Selvbevidsthed,
hvad saa? Saa er den gradvise Overgang sprængt, der
er skeet en Forvandling: Dyrene, hvorledes de for Resten
see ud, ere ikke Dyr, men Mennesker. Lovgivende Faar
og selvbevidste Høns ere Pjat; Heste med Aand maae
være mere end Heste; hvorledes skulde vel Heste kunne
blive Geologer?

Hvor uendelig dyb Forskjellen er mellem Dyrets og
Menneskets Sjæl, kan Enhver overbevise sig om, der vil
lægge nogen Vægt paa Karakter og Villie. Vel har Dyret
Karakter, desværre ikke sjældent en mere konsekvent
Karakter end et Menneske; men Grunden er let at opdage.
Dyrets Karakter er en Naturbestemmelse, der ligger
i dets Art; Dyret er fra Naturens Haand, hvad det
skal og bør være. Dersom det, at leve overensstemmende
med sin Natur, er Dyd, saa ere alle Dyr dydige; det er
mere, end man kan sige om Menneskene. De, som ikke
% --- p. 494 ---
\setcounter{footnote}{0}%
\opage{494}ret kunne faae Øie paa den egenlige Forskjel mellem
Menneske og Dyr, hvad det Positive angaaer, gjøre vel i
at begynde med det Negative. Thi Mennesket kan gjøre,
hvad Dyret ikke kan, enervere sig selv og forspilde sit
Liv; Mennesket kan hænge sig, skyde sig, »sprætte Maven
op paa sig« og i det Hele udvise, som Lærebogen udtrykker
det, »en Grumhed imod sig selv, hvorpaa der
ikke findes Exempler blandt de vildeste Dyr«; her er en
radikal Forskjel. Mennesket har af Naturen ingen Karakter,
men vel Karakteranlæg; dets eiendommelige Bestemmelse
er, at det med selvbevidst Villie skal udvikle
sine Anlæg og give sig selv sin Karakter: Karakteren er
Menneskets virkelige Selv, og det virkelige Selv er i Villien.

Der er en hel Kreds af Livsspørgsmaal, hvis Behandling
forkvakles i de almindelige Lærebøger, fordi deres
mulige Løsning udelukkende beroer paa Karakterdannelse.
Saaledes Spørgsmaalet om Villiens Frihed. De, der paastaae,
at Villien, uafhængig af alle Naturtilskyndelser, kan
vælge, beslutte og handle, som den vil, fare frem i Blinde.
Forholdt det sig saaledes, at Villien kunde frigjøre sig
fra alle Bevæggrunde, maatte det jo staae i Menneskets
Magt at omskabe sit Væsen fra Grunden af og gjøre sit
Sind og sine medfødte Anlæg til, hvad det skulde være.
Erfaringen taler for eftertrykkelig mod en saa vild Vilkaarlighedslære,
til at vi skulde spilde nogen Tid paa at
gjendrive den. Men at gaae til den modsatte Yderlighed
og gjøre Menneskets Villie til en naturbunden Yttring af
et medfødt Sindelag, som kun Opdragelse, Vane og Andres
Exempel i det Høieste kan tvinge ind under en vis
Form og lede i en vis Retning, er endnu farligere, thi i
denne Synsmaade er det noget Skuffende. At Villien
% --- p. 495 ---
\setcounter{footnote}{0}%
\opage{495}maa have en Naturgrund, og at denne Naturgrund kan
være forskjellig hos de forskjellige Individer, er umiskjendeligt.
Vi see det paa Børn. Et Barn paa fem
å sex Aar kan i visse Belysninger staae for den opmærksomme
Iagttager som en vel udrustet lille Person,
der er i Færd med at give sig en Karakter; et stridt
eller et blødt Sind, Heftighed, Sindighed, Dristighed, Frygtsomhed,
Aabenhed, Indesluttethed, kort, de forskjelligste
Tilsyneladelser snart af noget Ondt, snart igjen af noget
Herligt og Godt komme, ofte paa en høist overraskende
Maade, frem paa Sindets Overflade fra en dunkel Grund
og i mangfoldige Afskygninger. Da enhver aandelig
Virksomhed forudsætter en Naturbegavelse, og en Naturbegavelse
ofte kan yttre sig meget tidlig, saa er der
dem, der mene, at ogsaa Villien maa beroe paa en særegen
Naturgave, at det stærke Sind nødvendig maa give
en stærk Villie, det svage Sind en svag Villie, ja at man
allerede kan see paa Individet i en meget ung Alder, om
det nogensinde vil faae en Villie eller er bestemt til at
gaae Livet igjennem i Omstændighedernes Ledebaand.
Her have vi igjen en af disse halve Sandheder, som ved
at udgives for hele blive til Usandheder. Saalænge man
ensidig betragter Mennesket fra et zoologisk Standpunkt,
er det jo ganske sandt, at det forholder sig med den
driftagtige Villie som med det naturlige Mod. En har
Løvens Natur og er modig, en Anden Harens Natur og
er feig; En fik stærke Drifter, store Lidenskaber og blev
til en Mand af Villie, en Anden fik kun en døsig Drift,
et mat Sind og en lam Villie: det er, siger man, Naturbestemmelse
altsammen. Det høres ikke sjeldent, at den,
der bliver overbevist om en vitterlig Uret, uvilkaarlig
% --- p. 496 ---
\setcounter{footnote}{0}%
\opage{496}udbryder: »jeg vil slet ikke høre Noget; det er min Natur;
jeg er nu en Gang saaledes og kan aldrig blive anderledes«.
Der er i den Art Forsvar en egen Blanding
af Trods, Fortvivlelse og Undskyldning, som vidner om,
at Personen betragter sig selv zoologisk.

Den zoologiske Betragtning af Menneskets Villie staaer
i nøie Forbindelse med den statistiske. Ogsaa af Statistiken
udleder man Beviser mod Villiens Frihed, og Beviserne
ere saa meget kraftigere, som de ligefrem kunne
bygges paa Kjendsgjerninger og føres ved sikker Beregning.
Naar man har en Tabel over de Handlinger, som
igjennem en vis Aarrække have gjentaget sig, overensstemmende
med Befolkningens Karakter i et vist Land, saa
kan man med Bestemthed forudsige, hvad den frie Villie
hos denne Befolkning maa ville udrette i det kommende
Aar. Man kan ikke blot i Almindelighed forudsige, hvormange
Selvmord, hvormange Ægteskaber og hvormange
Tyverier, der ville indtræffe, men man kan gaae i det Enkelte
og udregne, hvad Selvmord angaaer, hvormange der
skulle hænge sig og hvormange der skulle drukne sig,
samt i hvilke Maaneder de frie Individer ville foretrække
Hængningen for Drukningen og omvendt; hvad Ægteskaberne
angaaer, hvormange der ville gifte sig med Enker
og hvorledes Antallet af dem, der gifte sig med gamle
Enker, vil forholde sig til Antallet af dem, der gifte sig
med unge; hvad Tyverierne angaaer, hvormange Marktyverier,
hvormange Indbrudstyverier og hvormange Lommetyverier,
den tyvagtige Villie i Aarets Løb vil udføre
o. s. fremdl. Men kan man i Statistiken forud beregne
Ægteskaber, Selvmord, Tyverier o. desl. med samme Sikkerhed,
hvormed man i Astronomien kan beregne Sol- og
% --- p. 497 ---
\setcounter{footnote}{0}%
\opage{497}Maaneformørkelser, saa synes jo heri at ligge et slaaende
Bevis for, at Menneskets Karakter i det Hele er naturbunden
som Dyrets, at den saakaldte menneskelige Frihed
bevæger sig efter lige saa nødvendige Love, som Kloderne
i Himmelrummet.

Til at blive hele dette Vilderede med Geologien,
Zoologien og Statistiken kvit gives der kun een Udvei:
bryde af, begynde forfra, tage Spørgsmaalet om Villiens
Frihed for, hvad det er: ikke et Videnskabsspørgsmaal,
men et Karakterspørgsmaal. Skal Forskjellen mellem
Menneske og Dyr erkjendes for mere end en Gradforskjel,
da er det ikke nok, at Mennesket opfatter sig som
tænkende og falder i Tanker over den mærkværdige Lighed,
der er imellem et Dyr og et Menneske. Mennesket
maa fra Grunden af begynde med at besinde sig paa sin
oprindelige Bestemmelse, tage sig sammen og søge sit
egenlige Selv i Villien. Hvad Tænkning er, lære vi ved
at tænke; men hvad Villie er, lære vi kun ved at ville.
At ville er Selvets inderste Handling, den Handling, hvorved
det hæver sig op af sin Naturgrund, sin Mulighed,
frembringer sig selv, gjør sig til noget Virkeligt, Eiendommeligt,
Selvstændigt. Denne Handling gaaer ikke
strax ud paa at faae Dette eller Hint udrettet i Verden,
den retter sig ikke mod det Udvortes, det Objektive, men
ubetinget mod det Indvortes, det Subjektive.

Der er i Nutiden to farlige Klipper, paa hvilke Individet
let kan strande, naar det for Alvor skal til at
bruge sin Villie og blive til Menneske, den ene er en
misforstaaet Praxis, den anden en misforstaaet Theori.
Den praktiske Misvisning er, at Villien uden dybere Selvbesindelse
vender sig udad og haster med at faae Noget
% --- p. 498 ---
\setcounter{footnote}{0}%
\opage{498}udrettet; Handlekraften offres paa det Endelige i timeligt
Erhverv og mange Slags borgerlige Forretninger. Saa
fremkommer der Noget, der ligner et Selv, en Karakter,
og dog er det hult inden i; thi dette Selv veed ikke,
hvad det selv er, om det virkelig er Noget i sig selv eller
blot et organisk Produkt af Muskler og Nerver. Det er
uhyggeligt at see, hvorledes et Menneske kan anstrenge
sig til det Yderste, for at vise Verden, hvor høit han
elsker sit eget Selv, og naar saa Alting lykkes ham, og
han med sit kjære Selv ret bilder sig ind at være ovenpaa,
saa har han dog intet Selv. Hvor blev det af? Det
blev borte for ham; det svandt hen i Forretninger, det
gik under i det Udvortes. Men er den praktiske Misvisning
farlig, da er den theoretiske endnu langt farligere.
Theorien vil, at det tænkende Væsen ikke skal forhaste
sig med det Udvortes, men fordybe sig i sig selv og besinde
sig paa, hvad Villie er, før det begynder at ville.
At ville er at tænke: den rene Villen er den rene Tænken.
Saa sidder da Tænkeren hen og tænker: først det
Objektive og saa det Subjektive, først Nødvendigheden og
saa Friheden, først det Zoologiske og saa det Psykologiske;
nu er Tanken i Metafysiken, nu i Ethiken, nu i Historien,
nu i Statistiken; og hvad kommer der ud af det
Hele? At det Subjektive er objektivt, at Frihed er Nødvendighed,
at Selvet er en almindelig Tænken. Theorien
er konsekvent; den ender som den begynder: den begynder
med Tænken og ender med Tænken. Den objektive,
almindelige Tænken er Alt, har Alt, tager Alt, drukner
Alt, Tænkeren, Selvet, Villien, i evig Nødvendighed;
saa er Villien tænkt, saa kan Tænkeren begynde at ville
--- tænke.

% --- p. 499 ---
\setcounter{footnote}{0}%
\opage{499}Først naar Frihedsspørgsmaalet forvandles til et Karakterspørgsmaal,
er Alting klart. Mennesket er ikke
Menneske i samme Forstand, som Dyret er Dyr; Mennesket
skal ved sin Villie gjøre sig til, hvad det væsenlig
er, til virkeligt Menneske. Hvad det virkelige Menneske
angaaer, er det ikke Ansigtet, de to Hænder og
den opreiste Gang, der gjør Udslaget; nei, det er Villien:
saadan en Villie, saadan et Menneske, saamegen Villie,
saameget Menneske, hverken Mere eller Mindre. Fra
dette Synspunkt er det da ogsaa let at komme til Rette
med Zoologien. Villiesanlæget sætter Forskjel, en for Tid
og Evighed gjældende Forskjel mellem Menneske og Dyr,
men ingen Forskjel mellem Menneske og Menneske, thi
hvert Menneske har Gaven. Sindskræfter og Lidenskaber
ere hos de forskjellige Mennesker høist forskjellige; men
det at ville, eller rettere: det at \emph{kunne} ville er fælles
for alle Mennesker. Selvbevidsthed, Selvbestemmelse
og Villie ere i Anlæget Et og det Samme. Her spørges
ikke om, hvor stor Villiens Naturkraft er, thi den driftagtige,
naturbundne Villie staaer unægtelig i Forhold til
individuel Begavelse; men her spørges om Selvbestemmelsens
indre Væsen, om Villiens Anstrengelse for at
arbeide sig op til virkelig Villie. Lad et Menneske være
nok saa skrøbeligt og faldet i nok saa mange Fristelser;
dersom et saadant Menneske med god Samvittighed kunde
sige til sig selv: jeg veed, at jeg ærlig har opbudt hele min
Villies Kraft og gjort Modstand til det Alleryderste, da
vilde dette Menneske midt i sin timelige Fornedrelse staae
som et ophøiet Selv med en Villie mere heroisk end
Alexanders, Cæsars eller Napoleons. Sagen er, naar det
gjælder om at ville gjøre, hvad Selvet erkjender for
% --- p. 500 ---
\setcounter{footnote}{0}%
\opage{500}ædelt og godt, eller at ville modstaae, hvad det erkjender
for lavt og slet, at da den menneskelige Villie aldrig vil
til det Yderste. Denne Svaghed, en ubegribelig Selvmodsigelse,
er fælles for Helten og det ubetydeligste Menneske.
Jo mere udviklet Villiesbevidstheden er hos et
Menneske, desto dybere føler han, at han har gjort sig
til en mere eller mindre forkvaklet Karakter, der bærer
paa et Ansvar; jo Mindre der er i et Individ af et virkeligt
Menneske, desto lettere stiller det sig tilfreds med
Undskyldninger; jo dybere udviklet Mennesket er, desto
Mere vil det have at bebreide sig selv. Selvbebreidelse,
dette strenge Mærkeskjel mellem Dyr og Menneske, er et
talende Vidnesbyrd om Frihedens Oprindelighed, om Villiens
indre Uendelighed. Fra dette Synspunkt er det let
at komme til Rette med Statistiken. Statistiken handler
jo slet ikke om, hvad Menneskene kunne gjøre, men hvad
de pleie at gjøre; den udpeger ikke de enkelte Individer
og paastaaer, at disse Enkelte, hver især, uundgaaelig
maae ville Dette eller Hint, tvertimod! om de Enkelte
som Enkelte veed Statistiken Intet. Den tager Individerne
mængdevis, betragter dem som Atomer i Massen
og kommer saa ganske rigtig til et Gjennemsnitsresultat,
som klarlig viser, at Individerne massevis, civiliserede som
ucivilicerede, endnu ikke ere hele, men halve Mennesker.

At Spørgsmaalet om Villiens Frihed hænger paa det
Nøieste sammen med Spørgsmaalet om Religion og Tro,
vil vist Ingen nægte. Er nu Frihedsspørgsmaalet, oprindelig
taget, ikke et Videnskabsspørgsmaal, men et Karakterspørgsmaal,
da maa det Samme jo gjælde om Troesspørgsmaalet.
»Viis mig din Tro af dine Gjerninger«,
denne Sætning, som, uagtet dens Sandhed er baade stor
% --- p. 501 ---
\setcounter{footnote}{0}%
\opage{501}og klar, har vakt saamange Stridigheder, vildledt saamange
Sjæle og forvirret saamange Hjerner, vil vel omsider
kunne befries for al Misforstaaelse og stilles i sit
rette Lys, naar det hedder: »\emph{Viis mig din Tro af din
Karakter}!« Karakter er noget Mere end en Sum af
Gjerninger, og dog høre Gjerninger med til en Speiling
af den virkelige Karakter; Gjerningerne ere som Lineamenter
i Karakterens Fysiognomi, men Karakteren selv
er Inderlighed, er Aand og Væsen. Har det udhulet Katholicismen
og slaaet Pietismen med aandelig Ufrugtbarhed
at ville hænge sig i Gjerninger, da har den lutherske
Orthodoxi tilvisse ogsaa maattet lide for at ville holde
Gjerningerne udenfor Troens Væsen. Begge Afveie have
deres Udgangspunkt i de theologiske Lovstridigheder, i
hvilke en skolastisk Skarpsindighed er blevet fristet til
at adskille det, som absolut maa være forenet i den religiøse
Karakter. Det er fuldkommen sandt, hvad Luther
har lært, at Troen er Væsen, at dette Væsen er Liv, at
Livet er i Inderlighed en Guds Kraft til Salighed; men
naar Inderlighedens Kræfter røre sig, saa komme Gjerningerne
frem, og saa have vi Karakteren. Udelades Karakterfordringen,
saa deler Bevidstheden sig almindeligvis
i to Halvdele; den ene Halvdel tilhører Evigheden, den
anden Timeligheden, og saa kan man igjen gaae forskjellige
Veie. Kommer det slet ikke an paa Gjerninger, men
kun paa Troen, da ligger det jo nær, at opfatte Troen
som en Sag for sig; et Sindelag, der antager den rene
Lære, lader sig opbygge af de i en Søndagskirke anordnede
Ceremonier, og ved Siden heraf har sine Gjerninger
i det Timelige, passer sit Embede, sin Næringsvei,
lever med al borgerlig Retsindighed i det Borgerlige
% --- p. 502 ---
\setcounter{footnote}{0}%
\opage{502}udenfor det Evige. Man har med Rette gjort gjældende, at
»de Kristne« nuomstunder ikke ere til at skjelne fra andre
skikkelige Mennesker. Hvilken Forskjel er der vel i
den dannede Verden mellem Jøder og Kristne? Jøderne
gaae i Synagogen om Løverdagen, de Kristne gaae i Kirke
om Søndagen; det er hele Forskjellen. Paa Høidepunktet
af den nuværende europæiske Civilisation gives der ---
takket være den klassiske Intelligens --- hverken Kristne eller
Jøder eller Hedninger: vi ere alle neutrale Kulturmennesker.
Men naar Intelligensen ikke tør opgive Religionen, uagtet
den fra Gjerningerne adskilte Tro er for inderlig mat til
at give sig virkelig Karakter i en virkelig Verden, saa er
Følgen heraf en aandelig Uddøethed, der er saa meget mere
ødelæggende for det indvortes Menneske som den ikke
sjældent behager sig selv i den Indbildning, at man dog
virkelig har Religion, naar man kraftig holder paa Statskirken
og nærer en dunkel Respekt for Dogmerne. Hvorvidt
en saadan Uddøethed kan gaae i en gammel Kultur,
og hvor let den lader sig forene med en indbildt Agtelse
for noget Helligt, Idealt og Ophøiet, derpaa er følgende
Exempel et talende Bevis.

Missionæren Hue fortæller\footnote[1]{\textit{L'empire chinois} 1, 173; jvfr. Oscar Ræder: »Historiske Love«. Kbhvn. 1871, S. 20---21.}, hvor overordenlig vanskeligt
det er at gjøre Proselyter for Kristendommen i
Kina, da Indifferentismen, den religiøse Karakterløshed,
er saa stor, at Intet kan overvinde den. Det er langt
fra, at Missionæren kun modtages med Vrede; tvertimod!
Kineseren hører paa den Fremmede med Fornøielse; det
interesserer ham og vækker hans Videbegjærlighed. Han
% --- p. 503 ---
\setcounter{footnote}{0}%
\opage{503}antager, billiger, alt det, der siges ham, og gjør ikke den
ringeste Vanskelighed eller Indvending; han finder, at
det, der siges, er sandt, smukt og fortræffeligt, ja han
optræder selv som Prædikant, og nu hører man ham tale
om at plyndre Gudebillederne til Gunst for Kristendommen.
Han beklager de Menneskers Forblindelse, der
knytte sig til denne Verdens Goder, og han vilde, om
det forlangtes, holde en prægtig Tale om den Lykke at
kjende den sande Gud, at tjene ham, og paa denne Maade
vinde det evige Liv. Idet man hører ham, skulde man
troe, at han allerede var Kristen, men han er ikke kommet
et Skridt videre. Og dog maa man ikke troe, at
hans Ord ikke til en vis Grad skulde være oprigtige; han
troer, hvad han siger, han taler gjerne derom, men bestandig
som om en Ting, der ikke er til for ham og ikke
vedkommer ham. Den kinesiske Beundrer af den kristne
Tro ender da stundom med at erklære Missionæren, der
beder ham at paaskynde sin Omvendelse, at han, om han
end ikke er blind for den kristne Religions Skjønhed og
Høihed, dog ikke seer, hvorfor man skulde forøge Livets
Plager og Bekymringer mere end nødvendigt ved at beskæftige
sig med et andet Liv end det nærværende og
stræbe ud over sin Evne. Det Legeme, som vi see og
berøre, koster os hver Dag den største Pleie og Omsorg,
skulde vi da ovenikjøbet beskæftige os med en Sjæl, som
vi ikke see, og et Liv, som vi ikke kjende. --- Stakkels
Kineser, han har Tro, men mangler blot et Par Gjerninger.
Havde han været født i en europæisk protestantisk
Statskirkes Moderskjød, saa var det med Gjerningerne
faldet som af sig selv; thi han har jo Gaverne.
Man tænke sig, at en saadan Mand havde studeret Theo%
% --- p. 504 ---
\setcounter{footnote}{0}%
\opage{504}logi, og lært at præke i Pastoralseminariet: han var blevet til noget Stort; han var vist blevet mere end Provst,
han var blevet Bisp og Sandhedsvidne.

Ved at lade Gjerningerne gaae paa det Timelige og
Troen paa det Evige kan man ogsaa komme ind paa en
anden Afvei, nærmest den pietistiske. Den Følelse paatrænger
sig, at man her i det Timelige dog virkelig maa
gjøre Noget for det Evige; men hvad skal man gjøre?
Det Timelige er timeligt og bliver aldrig Andet. Saa
maa man da begynde med at forsage det Timelige. Hvormeget
skal man nu forsage? Ikke alt Timeligt; thi dersom
man bogstavelig forsager alt Timeligt, saa døer man
strax; man tager sig selv af Dage. Saa deler man da
de timelige Ting og Gjerninger i to Klasser. Den første
Klasse omfatter de nødvendige og tilladelige Gjerninger,
saasom: at spise, drikke, sove, stifte Familie, drive borgerligt Erhverv o. s. v.; det heri indeholdte Timelige skal
altsaa ikke forsages. Til den anden Klasse høre ikke
blot saadanne Laster og Udyder, som Politiet straffer,
men ogsaa Handlinger, som det sædelige Samfund ellers
anseer for uskyldige og tilladelige, f. Ex. at dandse, spille
Kort og gaae paa Komedie, ja, hvad mere er: at elske
sit Folk og forsvare sit Fædreland, sørge ved dets Ulykker,
glæde sig ved dets Hæder, ære dets Minder og arbeide
med paa dets Fremtid, regnes iblandt de Daarskaber, af
hvilke Djævelen især fletter sine Garn, naar han lumsk gaaer
omkring og fanger Sjæle i det Timelige. Men hvo seer
nu ikke, at en registeragtig Opregnen, af hvad der skal
forsages og ikke forsages, mere egner sig til at udklække
forskruede Individer end til at optugte Karakterer. Vi
ære den pietistiske Alvor, naar vi advare mod dens Feil%
% --- p. 505 ---
\setcounter{footnote}{0}%
\opage{505}greb. Det er med Hensyn til de store, maaske forfærdelige,
Aandskriser, som Fremtiden stiller i Udsigt, en sand
Lykke for vort Fædreland, at det i sin Literatur eier en
Forfatter, som ud fra et pietistisk Standpunkt har paa
den mest omfattende og gjennemgribende Maade opløst
hele Pietismen; det er S. Kierkegaard.

At bryde med Verden, forsage Verden, afdøe fra Verden,
det er og bliver, saalænge der er Kristendom til, en
kristelig Grundfordring, en uafviselig Karakterfordring,
idet der da ved Verden forstaaes, hvad vi alle vide, det
sandselige, selviske, timelige, forfængelige Væsen. Fordringen
er negativ, en umiskjendelig ægte Lovens Fordring.
Hvori bestaaer nu det Pietistiske? Ikke deri, at
»Idealerne forkyndes«, at Fordringen gjøres gjældende i
hele sin Udstrækning; thi det er jo netop det Kristelige,
og Enhver, der under Mildhedens Maske paa Blødsødenhedens
Vegne tillader sig at fortie, fortone eller tilsløre
den hellige Fordring, er ganske vist ikke en kristelig
Karakter, endnu mindre et Sandhedsvidne, men en
Forræder, der kysser den Korsfæstede. Nei, det Pietistiske
bestaaer i noget Andet; det bestaaer deri, at
den negative Lovens Fordring stilles i Evangeliets Sted,
stilles op for sig, som om den selv indeholdt Betingelserne
for sin Opfyldelse. Det er dette, der er det Uevangeliske.
Nu laa det ikke blot i S. Kierkegaards exklusivt
nervøse Individualitet --- enten \emph{Lydighed} i Tro eller en
Plads i Daarekisten --- og i hans ganske eiendommelige
Karakterudvikling, men ogsaa i den Maade, hvorpaa han
opfattede sin Mission, sin Kamp imod hvad han ansaae
for en epikuræisk enerveret, i Verdslighed fortumlet, af
»Præstevrøvl« bedøvet Samtid, at stille Fordringen, som
% --- p. 506 ---
\setcounter{footnote}{0}%
\opage{506}han gjorde. Derved skete det, at hans Fremstilling af
Kristendommen, skjøndt den vitterlig indeholder alle Evangeliets
Grundbestemmelser, fik et eget Skær, et Drag af
Pietisme.

Men Pietisterne skulle neppe glæde sig, naar de betænke,
hvorledes deres Verdensforsagelse her er gjennemført.
At sætte den religiøse Alvor i at afholde sig fra
Skuespil, Kort og Dands, klynke over Verdens Ugudelighed
og for Resten være som Andre --- en Smule Tyv i
sin Næring, er ikke Alvor; naar en saadan Forsagelse vil
udgive sig for mere end, hvad den i mange Tilfælde vel
kan være, en tilraadelig Varsomhed imod at falde i Fristelse,
vil gjælde for fuld Alvor, hellig, dyb og gudelig
Alvor, da bliver jo den Slags gudelige Alvor, hvad Enhver
kan see, ligefrem af S. K. »udleveret til komisk Behandling«.
Nei, Fordringen er, at Alt, ubetinget Alt, ---
ogsaa Næringsveien --- skal forsages. Fordringen er, at
Timelighedens Interesser alle skulle opbrændes i Aandens
Ild, og det Selviske fortæres af Flammen. Det Ubetingede
i Fordringen, at Alt, alt Timeligt uden Undtagelse,
skal forsages, er det Ligelige, og denne Ligelighed er
igjen det i Alvoren Ophøiede.

Men lad os nu see, hvorledes den ophøiede Alvor
tager sig ud, naar Idealerne forkyndes, og en menneskelig
Villie skal prøve paa at sætte det Evige i Karakter.
»Overalt, hvor det Kristelige er,« siger S. Kierkegaard i
»Kjærlighedens Gjerninger« (S. 63), »er ogsaa Selvfornægtelse,
der er Kristendommens væsenlige Form\dots{} Men
Elskov og Venskab er netop Selvfølelsens Høieste, er
Jeget beruset i et andet Jeg. Jo fastere de tvende Jeger
slutte sig sammen til at blive eet Jeg, jo mere slutter
% --- p. 507 ---
\setcounter{footnote}{0}%
\opage{507}dette forenede Selv sig selvisk ud fra alle Andre. I Elskov
og Venskabs Høieste blive de Tvende virkeligen eet
Selv, eet Jeg. Dette er kun forklarligt, fordi der i Forkjærlighed
er indeholdt Naturbestemmelse (Drift --- Tilbøielighed)
og Selvkjærlighed, der selvisk kan forene Tvende
i et nyt selvisk Selv. Aandens Kjærlighed derimod borttager
af mit Selv al Naturbestemmelse og al Selvkjærlighed,
derfor kan Kjærlighed til Næsten ikke gjøre mig til
Eet med Næsten i et forenet Selv\dots{} Er Du, m. T.
(S. 67) maaskee, som det hedder en Dannet, nu vel, jeg
er ogsaa dannet; men dersom Du troer ved Hjælp af
»Dannelse« at komme dette Høieste nærmere, da tager
Du høiligen feil. Og her stikker just Vildfarelsen; thi
Dannelse ønske vi alle, og Dannelse har i Et væk det
Høieste i Munden, ja, ingen Fugl, der kun har lært eet
eneste Ord, raaber mere uafbrudt paa dette eneste Ord,
og Kragen ikke mere uafbrudt paa sit eget Navn end
Dannelsen altid paa »det Høieste«. Men det Kristelige
er ingenlunde Dannelsens »Høieste«, og det Kristelige
optugter just gjennem Forargelsens Frastød\dots{}. Du
\emph{skal} elske, Du skal elske din Næste som dig selv! (S.
70—71). Lad ikke være at troe, fordi Budet næsten forarger
Dig, fordi Talen ikke lyder indsmigrende som Digterens,
der med sine Sange indynder sig i din Lykke...
Nei, elsk den Elskede trofast og ømt, men lad Kjærlighed
til Næsten være det Helliggjørende i Eders Forenings
Pagt med Gud; elsk din Ven oprigtigt og hengivent, men
lad Kjærlighed til Næsten være, hvad I lære af hinanden
i Venskabets Fortrolighed med Gud!.. Som Kristendommens
glade Budskab er indeholdt i Læren om Menneskets
Slægtskab med Gud, saa er dens Opgave Menne%
% --- p. 508 ---
\setcounter{footnote}{0}%
\opage{508}skets Lighed med Gud\dots{}. Forsaavidt Du elsker den
Elskede, ligner Du ikke Gud, thi for Gud er ingen Forkjærlighed\dots{}.
Forsaavidt Du elsker din Ven, ligner
Du ikke Gud, thi for Gud er ingen Forskjel. Men naar
Du elsker Næsten, da ligner Du Gud.«

Enhver, der alvorlig vil lægge sig paa Sinde, hvorledes
Modsætningen mellem Selvkjærlighed, det Ukristelige,
og Aandskjærlighed, det Kristelige, her fremstilles,
maa nødvendig blive opmærksom paa, at Fremstillingen
indeholder to Forklaringer, der forklare hinanden saaledes,
at Lys med Lys giver Mørke.

Første Forklaring. Den Kristne har ingen Elsket,
ingen Hustru, ingen Ven; han har kun Næsten! Vil Nogen
gjøre ham opmærksom paa, at han nødvendig maa
staae i et andet Kjærlighedsforhold til sin Elskede, hvis
han er forlovet, end til andre Piger, til sin Hustru, hvis
han er gift, end til andre Koner; det hjælper ikke. Dersom
han virkelig er i Karakter af, hvad her sættes som
det Kristelige, vil han elske alle Piger, alle Mænd og
alle Koner i den hele Verden ligelig, i samme Aand, med
samme Inderlighed og paa samme Maade, thi han elsker
i dem alle --- Næsten. Vil Nogen minde ham om, at
han foruden den almindelige Kjærlighed til alle Mennesker
nødvendig maa nære en ganske særegen Kjærlighed
til visse Mennesker, at han vel i sin Elskede skal elske
sin Næste, men dog ogsaa omvendt i \emph{denne} Næste sin
Elskede, i \emph{denne} Næste sin Ven: det er forgjæves. Det
Øie, hvormed han skulde see den Elskede som den Elskede,
Vennen som Ven, er slaaet ud; det var jo et hedensk
Øje; saa ere de begge, den Elskede og Vennen
blevne borte for ham: paa deres Plads staaer Næsten.
% --- p. 509 ---
\setcounter{footnote}{0}%
\opage{509}Selv om han kunde opdage dem, hvorledes skulde han
vel kunne elske dem. Kjærlighed til den Elskede er Forkjærlighed,
Kjærlighed til Vennen er Forkjærlighed; Forkjærlighed
er Selvkjærlighed, Selvkjærlighed er selvisk, og
det Selviske ugudeligt.

Anden Forklaring. Endskjøndt den Kristne i alle
Mennesker kun elsker Næsten, saa har han dog Tilladelse
til at have en Ven og en Elsket, som han særlig tør
elske. »Elsk den Elskede trofast og ømt, men«....
»elsk din Ven oprigtigt og hengivent, men«\dots{}. Meningen
kan umulig være, at den Elskede og Vennen blot
skulle elskes som Næsten; i saa Fald vilde det gjentagne
Men jo være uden al Betydning. Tilladelsen, om man
vil, Formaningen maa forklares saaledes: Naar Du, den
Kristne, blot ikke glemmer at elske hvert Menneske uden
Forskjel, ogsaa din Ven, din Elskede, som din Næste, føi
kun saa lidt Forkjærlighed til: elsk --- med Forkjærlighed
--- din Elskede trofast og ømt; elsk din Ven med
oprigtig og hengiven Forkjærlighed! Men hvad var nu
Forkjærlighed ifølge den første Forklaring? Selviskhed,
Synd, Ugudelighed! »Forsaavidt Du elsker den Elskede,
ligner Du ikke Gud, thi for Gud er ingen Forkjærlighed.
Forsaavidt Du elsker din Ven, ligner Du ikke Gud, thi
for Gud er ingen Forskjel«. Og hvad er Forkjærlighed
ifølge den anden Forklaring? Lidt Aflad, som Forfatteren
tilbyder os, ikke for Penge, det var Synd at sige,
men gratis og af et godt Hjerte. Men er det dog ikke
--- fra Lovens Synspunkt --- som sagde Forfatteren til sin
alvorlige Læser med klare Ord: Nu veed Du, hvad Synd
er, at det er Forskjelskjærlighed, Forkjærlighed. Men
naar Du veed det, erkjender det, føler det og bøier Dig
% --- p. 510 ---
\setcounter{footnote}{0}%
\opage{510}dybt for Idealet; synd saa kun --- »trofast og ømt« med
den Elskede, »oprigtigt og hengivent« med din Ven.

At det ikke kan falde noget besindigt Menneske ind
at ville paadutte S. Kierkegaard en saa vanvittig Forvrængning
af det Hellige, behøver vel neppe at anmærkes;
heller ikke ligger Modsigelsen i Idealets Fordring
som Ideal. Nei, Grundmodsigelsen, den isnende, skærende
Modsigelse fremkommer derved, at Loven umiddelbart
skal kræve af den Troendes Villie at sættes i Karakter
som Lov. Det er dette, der med al gudelig Tale om
Gjenfødelse, Besegling og Helliggjørelse er det Pietistiske.
At udfolde Idealets Banner med Lovens Indskrift er at
opløse Pietismen.

Vi vende os nu til den Grundtvigske Anskuelse, for
at see, hvorledes det her forholder sig med Karakter og
Villie. Grundtvig begynder ikke med Loven, men med
Evangeliet; han taler ikke om Tro og Gjerninger, men
om Aandens Kraft i »det levende Ord«, om Livets Fornyelse
i Troen; han stiller ikke Aandens Krav i fjendlig
Modsætning til Menneskets Naturbestemmelse, men forliger
begge og finder Enhedens Udspring i Hjertet. Hvad
følger saa heraf? De store Bud i Loven: Du skal elske
Gud over alle Ting og din Næste som Dig selv, blive
herved hverken oversete eller tilbagetrængte, men fremstillede
i Evangeliets Lys, saa at det Befalende: Du \emph{skal},
forvandles til en Forjættelse: Det skal gives dig at \emph{kunne}.
Paa hvilken Side er der nu, naar det, vel at mærke, gjælder
ikke blot om Fordringen, men om Opfyldelsen, den
dybeste Alvor? Aabenbart paa Evangeliets. Loven siger:
Du \emph{skal}, Du \emph{skal}! og veed dog meget vel, at Mennesket
ikke kan, hvad det saaledes skal. Men at fordre
% --- p. 511 ---
\setcounter{footnote}{0}%
\opage{511}af et Væsen, at det absolut skal, hvad det vitterlig ikke
kan, er ensidig Alvor. Evangeliet siger: Du skal, thi Du
kan; det er mig, der giver Dig Kræfter. Just det, at
der ikke fordres Andet eller Mere af Mennesket, end hvad
Mennesket væsenlig kan, forhøier uendelig Fordringens
Alvor. Oprindelig taget, begynder Evangeliet ikke med
en Fordring, men med et Tilbud. Det fordrer ikke af
Mennesket, at han paa egen Haand skal opfylde denne
eller hin Lov, men det tilbyder en Gave, den Aandens
og Kjærlighedens Gave, hvorved det bliver Mennesket
muligt at opfylde Kjærlighedens Lov, endnu tydeligere:
hvorved det bliver den guddommelige Kjærlighed muligt
selv at fuldkomme sin egen Lov i Mennesket. Evangeliet
henvender sig ikke med et Pligtbud, en Forskrift, til
Menneskets Villie, begynder ikke med at sige: Du er forpligtet
til at troe paa den Korsfæstede, forpligtet til at
lade Dig døbe o. s. v. Nei, Evangeliet gaaer Inderlighedens
Vei, begynder med en Forjættelse og taler til
Menneskets Hjerte, thi hele Mennesket er i Hjertet. Gjør
Evangeliet saa maaske Troen til en dunkel Følelse, uden
Aand og klare Tanker, uden Karakter og Villie? Ingenlunde!
Men Evangeliet kjender ikke, vil ikke kjende nogen
endelig Udstykning af vort indre Væsen; det vil tvertimod
helbrede det Sønderrevne, hele det Delte og ophæve
Stykværket: Følelse og Tanke, Aand og Villie maae
i Sandhed være Eet, naar Mennesket skal komme til
Reisning og faae Karakter. Evangeliet kjender ikke, vil
ikke kjende nogen indre Splid mellem Aand og Natur;
derfor søger det strax uden Omsvøb ind i Hjertet, thi
kun i Hjertet er Aanden Natur, og Naturen Aand. En
naturfjendsk Aand og en aandsforladt Natur ere uhygge%
% --- p. 512 ---
\setcounter{footnote}{0}%
\opage{512}lige Forvrængninger af Liv og Sandhed. Vil saa Evangeliet
maaske lokke Menneskene til sig og vinde dem ved
at smigre det naturlige Hjerte? Ingenlunde! Evangeliet
kjender det menneskelige Hjerte fra Grunden af og veed
fuldt vel, hvorledes det staaer til derinde; men det veed
ogsaa, at Hjertets Vei er Kjærlighedens eneste mulige
Vei, at der er ingen anden. Dersom det kun gjaldt om
borgerlig Retskaffenhed, om Pligternes Opfyldelse i Lydighed,
kunde Aanden jo let bruge Magt, indjage Skræk
og fremtvinge Lydighed. Men naar Lydigheden skal gaae
i Eet med Friheden og have sin Grund i Kjærligheden,
kan Aanden ikke bruge Magt, thi Kjærlighed lader sig ikke
fremtvinge. Kjærlighed kan drive Frygten ud, men ingen
Frygt i Himlen eller paa Jorden vil være i Stand til at
drive Kjærligheden ind. Loven kan true, tvinge, tordne,
befale: Du skal, Du skal! til evig Tid; dens Trusel er
kold og gold og kan ikke føde Kjærligheden; men Evangeliet
føder Kjærligheden --- i Menneskets Hjerte. Dersom
der i en menneskelig Natur ikke er Andet end selvisk
Attraa, sandselig Begjær og Kjødets Lyst, da har
Kristus aldrig havt nogen menneskelig Natur, og hans
Evangelium er en Umulighed. Men naar Evangeliet gaaer
til Hjertet og gjør, i guddommelig Forstand, det Naturlige
aandeligt, og det Aandelige naturligt, da er det Alvor,
hellig Alvor; thi det er ikke Menneskets kjødelige
Hjerte, der raader for Evangeliet, men Evangeliets Aand,
der raader i Hjertet.

Paa dette Grundlag kunne »Idealerne forkyndes«,
Karakterer opbygges, Villien styrkes som Aandens Organ
og »Kjærlighedens Gjerninger« øves. Men det vil da
ogsaa sees, at hele Kjærlighedsbegrebet, opfattet fra den
% --- p. 513 ---
\setcounter{footnote}{0}%
\opage{513}menneskelige Side, maa undergaae en Forvandling. Selvkjærlighed
og Forkjærlighed ere radikale Bestemmelser,
uadskillelige fra al den menneskelige Kjærlighed, der fra
nogen Side, i nogen Mening, paa nogen Maade skal kunne
siges at have en naturlig Oprindelse. Rykke vi Selvkjærlighedens
og Forkjærlighedens Forgreninger op med
Rode, da gaaer Evangeliet bort og lader en naturfjendsk
Aand med en naturløs Begyndelse tilbage. Selvkjærlighed
og Forkjærlighed ere tvetydige Ord; Tvetydigheden viser,
at de kunne tilhøre baade det Laveste og det Høieste.
Tag Forkjærligheden bort, hvad bliver der saa af Moderkjærligheden?
Tag Selvkjærligheden bort, hvad bliver
der saa af Kjærlighed til Næsten? »Elsk din Næste som Dig
selv!« Men at elske sig selv er jo Selvkjærlighed. I selvisk
Forstand er Selvkjærlighed Egenkjærlighed, og Egenkjærlighed
uforenelig med Aandskjærlighed. For at magte
Egoismen og udrense Egenkjærligheden, indgaaer Evangeliets
Aand et naturligt Forbund med den dybere Selvfølelse,
med Selvets Trang til evigt Liv, altsaa med den
sande Selvkjærlighed; saa bliver Aandsfordringen fra Grunden
af en Naturtilskyndelse, der vækker Villien ved Selvbekymring,
og den ved Selvbekymring vakte Villie bliver
Aandens Organ. At Selvopholdelsesdriften er en stærk
Drift, det gjælder ikke blot i sandselig, men ogsaa i
aandelig Forstand. Er det først i Alvor blevet et Menneske
indlysende, at Egenkjærligheden er hans værste
Fjende, da vil han hade Egenkjærligheden, ikke fordi han
hader sig selv, nei fordi han elsker sig selv. At Aandsvillien
hader Egenkjærligheden og hader den grundig, just
fordi den hader af Selvkjærlighed, vil jo med andre Ord sige,
at Villien af al Kraft bekæmper det Selviske, hører ikke
% --- p. 514 ---
\setcounter{footnote}{0}%
\opage{514}op dermed, slaaer sig ikke til Ro, før Selvsygens Braad
er døvet, og Dødsfjenden kvalt. Det er et Aandens
Værk, og dog gaaer det naturlig til: i den sande Selvkjærlighed
er det Naturlige aandeligt, det Aandelige naturligt.
Med Forkjærligheden forholder det sig paa samme
Maade. Al Forkjærlighed er naturlig og maa, forsaavidt
den holder sig i sandselig Selviskhed, nødvendig forkastes
af Aanden. Men det er kun det selvisk Indsnævrende,
Partiskheden, Hildetheden, Forgabetheden i Forkjærligheden,
ikke Forkjærligheden selv, som Aanden hader. Kjærlighed
til Andre er for Mennesker umulig uden Forkjærlighed.
Den virkelige Kjærlighed maa, som Alt, hvad
der i Verden skal blive virkeligt, nødvendig begrændses.
En ubegrændset Kjærlighed, en fuldkommen ligelig, almindelig
Kjærlighed til alle Mennesker i al Almindelighed
er vel en Sindets Udvidelse i Fantasi, men endnu ingen
virkelig Kjærlighed, thi den er ikke Kjærlighed til noget
virkeligt Menneske. Skal Kjærligheden blive virkelig,
maa den begrændse sig; men idet den begrændser sig og
vælger sin Gjenstand, bliver den Forkjærlighed. Det forstaaes
af sig selv, at den falske Forkjærlighed maa forholde
sig til den sande Forkjærlighed, ligesom Egenkjærligheden
forholder sig til den sande Selvkjærlighed. Manden
skal, efter Pauli Formaning, elske sin Hustru, som
han elsker sig selv. Dersom Mandens Selvkjærlighed nu
ikke er Andet end Egenkjærlighed, saa kan hans Kjærlighed
til Hustruen, hvor inderlig, øm og sværmerisk den
for Resten tager sig ud, heller ikke være Andet end en
i Forkjærlighed omsat Egenkjærlighed, en Forkjærlighed,
der selvisk udelukker Kjærlighed til Næsten. Dersom
hans Selvkjærlighed derimod i Sandhed er en for Egoisme
% --- p. 515 ---
\setcounter{footnote}{0}%
\opage{515}luttret Aandskjærlighed, da vil hans Kjærlighed til Hustruen
ogsaa fra første Begyndelse af være Kjærlighed til
Næsten. Lad ham kun begynde med at have denne Næste
som »den Elskede« og være nok saa lyksalig forelsket;
dersom blot hans Villie er i Karakter af Tro, vil Evangeliets
Aand indestaae for, at den Ild, der brænder i hans
Sjæl, ufeilbarlig skal luttre det sandselige Væsen; lad kun
de første umiddelbare Naturtilskyndelser svinde bort med
Aarene, Kjærligheden vil dog ikke kølnes, men Inderligheden
voxe med foryngende Kræfter og give Karakteren
et Udtryk, der, endskjøndt det ikke er i Ungdommelighedens
Blomstring, dog derfor er lige saa rigt, ja rigere,
fordi det er aandelig forklaret og dog derhos ikke mindre
naturligt. Har nogensinde nogen Mand bevaret og næret
en saadan Kjærlighed til sin Hustru, da behøves her ingen
Poesiens Herolder for at udbasune hans Erotik: Religionen
vedkjender sig ham, at hans Liv har været i
Aandens Tjeneste; Evangeliet giver ham det Vidnesbyrd,
at han har elsket sin Næste som sig selv. Vil Nogen
indvende, at en Begrændsning, som den her antydede,
staaer altfor lavt under Idealets Fordring, at det ikke er
nok at elske et eneste Menneske, allermindst sin Hustru,
men at man skal elske ogsaa fremmede Mennesker, mange
Mennesker, ja alle Mennesker: da maa hertil svares, at
det netop i denne Sag ikke kommer an paa et Flertal.
Næsten har, mærkværdigt nok, intet Flertal. Vi have
ikke mange Næster, men kun en Næste, hver Gang i
hvert virkelig personligt Forhold Næsten, \emph{denne} vor Næste.
Næsten er ikke blot Een i Alle, men ogsaa Alle i
Een; den, der i Sandhed elsker et eneste Menneske som sin
Næste, han elsker af Hjertet alle Mennesker, hvert Menneske
% --- p. 516 ---
\setcounter{footnote}{0}%
\opage{516}uden Undtagelse, thi dersom han i sit Hjerte gjør Undtagelse,
saa er hans Kjærlighed lige saa vel til de Flere
som til den Eneste et egoistisk Bedrag. Han kan for
den Sags Skyld gjerne elske ti Mennesker, dersom han
gjør en Undtagelse imod det ellevte, saa er det Hele en
Indbildning. Ligeligheden sigter paa Muligheden, thi den
angaaer Sindelaget; men Begrændsningen sigter paa Virkeligheden,
thi den angaaer Handlingen. Den, der har
lært at elske sin Næste som sig selv, har meget i Verden
at elske. Han elsker ikke blot sine Nærmeste, sin
Familie, sine Venner, men ogsaa sit Folk, sit Fædreland.
At ville elske sin Næste uden at elske sin Familie er
maskeret Egenkjærlighed; at ville elske sin Familie og
sine Venner, uden at elske sit Folk og sit Fædreland, er
falsk Forkjærlighed. At elske sit Folk og sit Fædreland
med sand Forkjærlighed er uendelig forskjelligt fra at
være forgabet i sin egen Stamme og forgude saadanne folkelige
Dyder, som i al deres Naturlighed dog altid ere parrede
med tilsvarende Unoder og Lyder; at elske sit Folk
med Forkjærlighed kan, saa sandt Forkjærligheden er i
Karakter af Aand, umulig lede til enten at hade eller
ringeagte eller miskjende det Gode og Fortræffelige, der
findes hos andre Folk, da en inderlig stærk og levende
Forkjærlighed for sit eget Folk tvertimod er den eneste
mulige Vei til at elske det Menneskelige i virkelige Mennesker,
den eneste Mulighed til at Kjærligheden atter kan
blive Kjød og Blod, at der kan komme Sundhed, Kraft
og Marv i den ellers saa blege, almindelige Menneskekjærlighed.

\emph{Viis mig din Tro af din Karakter}! Det Timelige
skal forsages, og dog skal det kraftigen bruges; den
% --- p. 517 ---
\setcounter{footnote}{0}%
\opage{517}Kristne skal »afdøe fra Verden«, og dog skal han, just
han, leve i Verden, virke i Verden, arbeide paa at føre
Menneskelivet fremad til rigere, sandere, fyldigere Udvikling
i Verden. »Idealerne skulle forkyndes«, de høie Forbilleder
paa Lidelse, Opoffrelse, Sjælekamp og Forsagelse
stilles frem i det Herlighedens Lys, der straaler med evig
Glands, og Lovens optugtende Bud: Du skal, Du skal!
bestandig lyde med al sin Strenghed; og dog skulle vi
Mennesker uden Skræk for Loven, i Frihedens, Selvforstaaelsens,
Evangeliets, Kjærlighedens Aand begynde ---
ret menneskeligt med det Guddommelige, ret naturligt
med det Aandelige, skulle i al Frimodighed med Glæde
og Haab paa det Uforkrænkelige arbeide og virke midt
i al denne Forkrænkelighed, denne Timelighed, denne
Verdslighed: det er for et overfladisk Syn modsigende
Fordringer. Viis mig din Tro af din Karakter! thi disse
Fordringer ere Opgaver, paatrængende Opgaver, ikke for
objektiv Viden og ræsonnerende Kritik, men for Karakter
og Villie.
\signature{R. \emph{Nielsen}.}

\end{document}
